% Gemini theme
% https://github.com/anishathalye/gemini

\documentclass[final]{beamer}

% ====================
% Packages
% ====================

\usepackage[T1]{fontenc}
\usepackage{lmodern}
%\usepackage[size=custom,width=118.9,height=84.1,scale=1.0]{beamerposter}
\usepackage[orientation=portrait,size=a0,scale=1.05]{beamerposter}
\usetheme{gemini}
\usecolortheme{gemini}
\usepackage{graphicx}
\usepackage{booktabs}
\usepackage{tikz}
\usepackage{pgfplots}
\pgfplotsset{compat=1.14}
\usepackage{anyfontsize}
\usepackage{xcolor}



\usepackage{todonotes}
\usepackage{algorithm}
\usepackage[noend]{algpseudocode}
\usepackage{xspace}

%\usepackage{amsthm}
%\newtheorem{lemma}{Lemma}
%\newtheorem{theorem}{Theorem}

\usepackage{cleveref}
\Crefname{definition}{Definition}{Definitions}
\crefname{definition}{Def.}{Defs.}
\Crefname{theorem}{Theorem}{Theorems}
\crefname{theorem}{Theorem}{Theorems}
\Crefname{lemma}{Lemma}{Lemmas}
\crefname{lemma}{Lemma}{Lemmas}
\Crefname{section}{Section}{Sections}
\crefname{section}{\S}{\S}
\crefname{algorithm}{Alg.}{Algs.}
\Crefname{algorithm}{Algorithm}{Algorithms}
\crefname{figure}{Fig.}{Figs.}
\Crefname{figure}{Figure}{Figures}
\crefname{table}{Tab.}{Tabs.}
\Crefname{table}{Table}{Tables}



% ====================
% Lengths
% ====================

% If you have N columns, choose \sepwidth and \colwidth such that
% (N+1)*\sepwidth + N*\colwidth = \paperwidth
\newlength{\sepwidth}
\newlength{\colwidth}
\setlength{\sepwidth}{0.025\paperwidth}
\setlength{\colwidth}{0.465\paperwidth}

\newcommand{\separatorcolumn}{\begin{column}{\sepwidth}\end{column}}








% vector
\newcommand{\V}[2][]{{\bm{#1\mathbf{\MakeLowercase{#2}}}}}
% matrix
\newcommand{\M}[2][]{{\bm{#1\mathbf{\MakeUppercase{#2}}}}}
% tensor
\newcommand{\T}[2][]{\boldsymbol{#1\mathscr{\MakeUppercase{#2}}}}


\newcommand{\Real}{\mathbb{R}}
\newcommand{\vc}[1]{\bm{#1}}
\newcommand{\frobenius}[1]{\left\| #1 \right\|}

\newcommand{\arraymul}{\emph{ternary multiplication}\xspace}
\newcommand{\pluseq}{\mathrel{+}=}
\newcommand{\Tra}{{\sf T}} 

% ====================
% Title
% ====================

\title{Minimizing Communication for Parallel Symmetric Matricized-Tensor Times Khatri-Rao Products}

\author{Grey Ballard \inst{1} \and Alimzhan Gumran \inst{2} \and \underline{Suraj Kumar} \inst{2} \and Kathryn Rouse \inst{3}}

\institute[shortinst]{\inst{1} Wake Forest University, USA \samelineand \inst{2} Inria and ENS Lyon, France \samelineand \inst{3} Inmar Intelligence, USA}



%Ballard, Grey (2); Gumran, Alimzhan (1); Kumar, Suraj (1); Rouse, Kathryn (3)
%Organisation(s): 1: Inria and ENS Lyon, France; 2: Wake Forest University; 3: Inmar Intelligence
% ====================
% Footer (optional)
% ====================

\footercontent{
	Suraj Kumar (Inria \& ENS Lyon, France) \hfill
%  \href{https://www.example.com}{https://www.example.com} \hfill
%  ABC Conference 2025, New York --- XYZ-1234 \hfill
  \href{mailto:suraj.kumar@inria.fr}{suraj.kumar@inria.fr}
}
% (can be left out to remove footer)

% ====================
% Logo (optional)
% ====================

% use this to include logos on the left and/or right side of the header:
% \logoright{\includegraphics[height=7cm]{logo1.pdf}}
% \logoleft{\includegraphics[height=7cm]{logo2.pdf}}

% ====================
% Body
% ====================

\begin{document}

\begin{frame}[t]
\begin{columns}[t]
\separatorcolumn



\begin{column}{\colwidth}
\begin{block}{Symmetric MTTKRP}
	An $R$-rank Canonical Polyadic (CP) decomposition of a $3$-dimensional symmetric tensor $\T{X}$ expresses it as a sum of $R$ symmetric rank-one tensors. More precisely, there is a matrix $\M{U}$ whose $r$th column is denoted by $\V{u}_r$, such that $\T{X} \approx \sum_{r=1}^R \V{u}_r \circ \V{u}_{r} \circ \V{u}_{r}$, where $\circ$ denotes an outer product. As $\T{X}$ is symmetric, therefore $\T{X}_{ijk}$ is same for all permutations of $i$, $j$ and $k$.
%	 \Cref{alg:cpals} specifies a method to compute symmetric CP decomposition of a $3$-dimensional tensor.
\begin{algorithm}[H]
	\caption{CP-ALS method to compute symmetric CP decomposition\label{alg:cpals}}
	\begin{algorithmic}[1]
		\Require $\T{X}\in \Real^{n\times n\times n}$ is a symmetric 3-dimensional tensor, initial factor matrix $\M{U}\in \Real^{n\times R}$
		\Ensure $\big(\lambda \in \Real^R, \M{U} \in \Real^{n\times R}\big)$ : A rank-$R$ CP decomposition of $\T{X}$
		\State Initialize each column of $\M{U}$ to a random unit vector (or based on a heuristic)
		\Repeat
		\State $\M{G} = (\M{U}^T\M{U}) \ast (\M{U}^T\M{U})$ \Comment{$\ast$ denotes the elementwise product}
		\State $\M{V} = \M{X_{(1)}}(\M{U}\odot \M{U})$
		\State $\M{U} = \M{V}\M{G}^\dagger$ \Comment $G^\dagger$ is the pseudo inverse of $G$
		\State Normalize columns of $\M{U}$ and store norms in $\lambda$
		\Until{convergence or the maximum number of iterations}
	\end{algorithmic}
\end{algorithm}
 $\M{X_{(1)}}(U\odot U)$ is the bottleneck computation of this algorithm and known as the symmetric Matricized-Tensor Times Khatri-Rao Product (MTTKRP). Here $\M{X_{(1)}}$ denotes the first unfolding matrix of $\T{X}$ and $\odot$ denotes Khatri-Rao product.

\end{block}
\begin{block}{Symmetric MTTKRP pseudocode}
		\begin{algorithmic}
			\Require $\T{X}\in \Real^{n\times n\times n}$ is a symmetric tensor and $\M{U}\in \Real^{n\times R}$ is a matrix.
			\Ensure $\M{V}\in \Real^{n\times R}$ such that  $\M{V} = \M{X_{(1)}} (\M{U}\odot \M{U})$.
			\State Initialize matrix $\M{V}$ to $0$
			\For{$i=1 \text{ to } n$}
			\For{$j=1 \text{ to } i$}
			\For{$k=1 \text{ to } j$}
			\Statex\Comment{Make all possible uses of $\T{X}_{ijk}$, $i\geq j\geq k$ }
			\For{$r=1 \text{ to } R$}

%			\State $\M{V}_{ir} \pluseq \T{X}_{ijk}\cdot \M{U}_{jr} \cdot \M{U}_{kr}$
			\If{$i\neq j$ and $j\neq k$}\Comment{An off-diagonal entry}
			\State $\M{V}_{ir} \pluseq 2\T{X}_{ijk}\cdot \M{U}_{jr} \cdot \M{U}_{kr}$
			\State $\M{V}_{jr} \pluseq 2\T{X}_{ijk}\cdot \M{U}_{ir} \cdot \M{U}_{kr}$
			\State $\M{V}_{kr} \pluseq 2\T{X}_{ijk}\cdot \M{U}_{ir} \cdot \M{U}_{jr}$
			\ElsIf{$i==j$ and $j\neq k$}\Comment{A non-central diagonal entry}
			\State $\M{V}_{ir} \pluseq 2\T{X}_{ijk}\cdot \M{U}_{jr} \cdot \M{U}_{kr}$
			\State $\M{V}_{kr} \pluseq \T{X}_{ijk}\cdot \M{U}_{ir} \cdot \M{U}_{jr}$
			\ElsIf{$i\neq j$ and $j==k$}\Comment{A non-central diagonal entry}
			\State $\M{V}_{ir} \pluseq \T{X}_{ijk}\cdot \M{U}_{jr} \cdot \M{U}_{kr}$
			\State $\M{V}_{jr} \pluseq 2\T{X}_{ijk}\cdot \M{U}_{ir} \cdot \M{U}_{kr}$
			\Else\Comment{A central diagonal entry, $i==j==k$ case}
			\State $\M{V}_{ir} \pluseq \T{X}_{ijk}\cdot \M{U}_{jr} \cdot \M{U}_{kr}$
			\EndIf
			\EndFor
			\EndFor
			\EndFor
			\EndFor
		\end{algorithmic}
%	\end{algorithm}
%\end{column}
	The $\M{V} = \M{X_{(1)}} (\M{U}\odot \M{U})$ computation can be computed directly element-wise as $\M{V}_{ir} =\sum_{j}\sum_{k} \T{X}_{ijk}\cdot \M{U}_{jr} \cdot \M{U}_{kr}$. We call the operation $\T{X}_{ijk}\cdot \M{U}_{jr} \cdot \M{U}_{kr}$ as \arraymul. The above algorithm takes advantage of symmetry and performs roughly half of the \arraymul operations compared to the direct way.

\end{block}

\begin{block}{Existing work on symmetric matrix computtions}
	\begin{minipage}{0.725\linewidth}
		\begin{itemize}
			\item To work with $a\times a$ \textcolor{red}{square block} of the symmetric matrix
			\begin{itemize}
				\item requires $2a$ elements of a non-symmetric matrix
			\end{itemize}
			\item For the same amount of computation with a \textcolor{orange}{triangle block}
			\begin{itemize}
				\item requires only $\sqrt{2}a$ elements of a non-symmetric matrix
			\end{itemize}
		\end{itemize}
	\end{minipage}\hfill
\begin{minipage}{0.265\linewidth}
	\begin{tikzpicture}[scale=0.75]

% Lower triangular matrix region
\fill[cyan!20] (0,0) -- (0,16) -- (16,0) -- cycle;

\fill[orange] (4,7.8) -- ++(4.2,0) -- ++(-4.2, 4.2) --cycle;

\fill[red] (1,1) -- ++(3,0) -- ++(0,3) -- ++(-3,0) --cycle;


\end{tikzpicture}
\end{minipage}

	


\textcolor{brown}{Approach to minimize communication}:  express computations in terms of triangle blocks

\heading{An example for $13\times13$ blocks on $13$ processors}


\begin{minipage}{0.495\linewidth}
	Each color represents the blocks handled by a \newline specific processor.
	\vspace*{2cm}
	\begin{itemize}
		\item Triangle block set of processor $7$: $\{2,7,9,13\}$
		\begin{itemize}
			\item Performs all computations associated with $6$ off-diagonal blocks: $(7,2), (9,2), (13,2), (9,7), (13,7), (13,9)$
		\end{itemize}
		\item Also performs computations associated with one diagonal block: $(7,7)$
	\end{itemize}
\end{minipage}\hfill
\begin{minipage}{0.475\linewidth}
\vspace*{-0.375cm}\begin{center}	
	\begin{tikzpicture}[scale=1.325]
	
	% colors
	\definecolor{c1}{RGB}{230,230,255}
	\definecolor{c2}{RGB}{210,235,255}
	\definecolor{c3}{RGB}{255,235,190}
	\definecolor{c4}{RGB}{220,220,220}
	\definecolor{c5}{RGB}{255,255,180}
	\definecolor{c6}{RGB}{255,210,220}
	\definecolor{c7}{RGB}{220,200,255}
	\definecolor{c8}{RGB}{255,220,190}
	\definecolor{c9}{RGB}{220,255,200}
	\definecolor{c10}{RGB}{190,255,210}
	\definecolor{c11}{RGB}{200,255,255}
	\definecolor{c12}{RGB}{255,210,255}
	\definecolor{c13}{RGB}{255,230,120}
	
	% matrix size
	\def\n{13}
	
	% grid
	\foreach \i in {0,...,13}{
		\draw (\i,0)--(\i,13);
		\draw (0,\i)--(13,\i);
	}
	
	% labels
	\foreach \i in {1,...,13}{
		\node at (\i-0.5,13.4) {\bfseries \i};
		\node at (-0.4,13.5-\i) {\bfseries \i};
	}
	
	% helper:
	% \cell{row}{col}{block}
	\newcommand{\cell}[3]{
		\fill[c#3] (#2-1,13-#1) rectangle (#2,14-#1);
		\node at (#2-0.5,13.5-#1) {\small #3};
	}



	

	
	%%%%%%%%%%%%%%%%%%%%%%%%%%%%%%%%%%%%%%%%%%%%%%%%%%
	% Lower triangular entries
	%%%%%%%%%%%%%%%%%%%%%%%%%%%%%%%%%%%%%%%%%%%%%%%%%%
	
	\cell{2}{1}{1}
	
	\cell{3}{1}{1}
	\cell{3}{2}{1}
	
	\cell{4}{1}{1}
	\cell{4}{2}{1}
	\cell{4}{3}{1}
	
	\cell{5}{1}{2}
	\cell{5}{2}{5}
	\cell{5}{3}{8}
	\cell{5}{4}{11}
	
	\cell{6}{1}{2}
	\cell{6}{2}{6}
	\cell{6}{3}{9}
	\cell{6}{4}{12}
	\cell{6}{5}{2}
	
	\cell{7}{1}{2}
	\cell{7}{2}{7}
	\cell{7}{3}{10}
	\cell{7}{4}{13}
	\cell{7}{5}{2}
	\cell{7}{6}{2}
	
	\cell{8}{1}{3}
	\cell{8}{2}{5}
	\cell{8}{3}{10}
	\cell{8}{4}{12}
	\cell{8}{5}{5}
	\cell{8}{6}{12}
	\cell{8}{7}{10}
	
	\cell{9}{1}{3}
	\cell{9}{2}{7}
	\cell{9}{3}{9}
	\cell{9}{4}{11}
	\cell{9}{5}{11}
	\cell{9}{6}{9}
	\cell{9}{7}{7}
	\cell{9}{8}{3}
	
	\cell{10}{1}{3}
	\cell{10}{2}{6}
	\cell{10}{3}{8}
	\cell{10}{4}{13}
	\cell{10}{5}{8}
	\cell{10}{6}{6}
	\cell{10}{7}{13}
	\cell{10}{8}{3}
	\cell{10}{9}{3}
	
	\cell{11}{1}{4}
	\cell{11}{2}{5}
	\cell{11}{3}{9}
	\cell{11}{4}{13}
	\cell{11}{5}{5}
	\cell{11}{6}{9}
	\cell{11}{7}{13}
	\cell{11}{8}{5}
	\cell{11}{9}{9}
	\cell{11}{10}{13}
	
	\cell{12}{1}{4}
	\cell{12}{2}{6}
	\cell{12}{3}{10}
	\cell{12}{4}{11}
	\cell{12}{5}{11}
	\cell{12}{6}{6}
	\cell{12}{7}{10}
	\cell{12}{8}{10}
	\cell{12}{9}{11}
	\cell{12}{10}{6}
	\cell{12}{11}{4}
	
	\cell{13}{1}{4}
	\cell{13}{2}{7}
	\cell{13}{3}{8}
	\cell{13}{4}{12}
	\cell{13}{5}{8}
	\cell{13}{6}{12}
	\cell{13}{7}{7}
	\cell{13}{8}{12}
	\cell{13}{9}{7}
	\cell{13}{10}{8}
	\cell{13}{11}{4}
	\cell{13}{12}{4}
	
	
	
	%%%%%%%%%%%%%%%%%%%%%%%%%%%%%%%%%%%%%%%%%%%%%%%%%%%%%%%
	% Diagonal tile entries
	%%%%%%%%%%%%%%%%%%%%%%%%%%%%%%%%%%%%%%%%%%%%%%%%%%%%%%%
	
	\cell{1}{1}{3}
	\cell{2}{2}{5}
	\cell{3}{3}{10}
	\cell{4}{4}{1}
	\cell{5}{5}{11}
	\cell{6}{6}{2}
	\cell{7}{7}{7}
	\cell{8}{8}{13}
	\cell{9}{9}{6}
	\cell{10}{10}{12}
	\cell{11}{11}{9}
	\cell{12}{12}{8}
	\cell{13}{13}{4}
	
	
	\end{tikzpicture}
	
\end{center}
\end{minipage}	

\begin{itemize}
	\item A triangle block partition for $n \times n$ blocks symmetric matrix exists on $P$ processors when $n=P=q^2+q+1$ for some prime $q$ (by the properties of projective planes or Steiner systems) 
%	\item Existence of the bound
\end{itemize}
\heading{References}
{\small
\begin{itemize}
	\item Beaumont et al. (2022) I/O-Optimal Algorithms for Symmetric Linear Algebra Kernels. In Proceedings of the 34th ACM Symposium on Parallelism in Algorithms and Architectures (SPAA '22)
	\item Al Daas et al. (2025) Communication Lower Bounds and Optimal Algorithms for Symmetric Matrix Computations. ACM Transactions on Parallel Computing (TOPC '25)
\end{itemize}
}

\end{block}



\end{column}

\separatorcolumn

\begin{column}{\colwidth}

\begin{block}{Communication lower bound}
	We obtain the following result by fixing the number of dimensions in \emph{Lemma 4.1} of (Ballard et al., Communication Lower Bounds for Matricized Tensor Times Khatri-Rao Product, (IPDPS '18)).
	
%	\begin{lemma}
\emph{
		Let $V$ be a finite set of points in $\mathbb{Z}^4$. Let $\phi_{ijk}(V)$ be the projection of $V$ on $i-j-k$ space, i.e., all points $(i,j,k)$ such that there exists a $r$ so that $(i,j,k,r) \in V$. Let $\phi_{ir}(V)$ be the projection of $V$ on $i-r$ plane, i.e., all points $(i,r)$ such that there exists a $(j,k)$ so that $(i,j,k,r) \in V$.  Define $\phi_{jr}(V)$ and $\phi_{kr}(V)$ similarly. 
		Then 
		$$|V| \leq |\phi_{ijk}(V)|^{2/3} \cdot |\phi_{ir}(V)|^{1/3}\cdot |\phi_{jr}(V)|^{1/3}\cdot |\phi_{kr}(V)|^{1/3}.$$
	}
%	\end{lemma}
We extend the above result for a symmetric set of points corresponding to the first three coordinates. Let $V$ be a finite set of points in  $\{(i, j, k, r)\in\mathbb{Z}^4 \, | \, i>j>k\}$, then
	$$6|V| \leq (6|\phi_{ijk}(V)|)^{2/3} \cdot |\phi_{ir}(V) \cup \phi_{jr}(V) \cup \phi_{kr}(V)|.$$
We consider that each processor performs equal amount of \arraymul operations, We solve the following optimization problem to determine the minimum number of entries accessed on a processor to perform its share of operations.
\heading{Optimization problem}
	$$\min |\phi_{ijk}(V)| + 2|\phi_{ir}(V) \cup \phi_{jr}(V) \cup \phi_{kr}(V)| \quad \text{subject to}\qquad\qquad\qquad\qquad\qquad\qquad$$
$$(6|\phi_{ijk}(V)|)^{2/3} \cdot |\phi_{ir}(V) \cup \phi_{jr}(V) \cup \phi_{kr}(V)| \geq 6\left(\frac{n(n-1)(n-2)R}{6P}\right) $$
%$$|\phi_{ijk}(V)|^2|\phi_{ir}(V) \cup \phi_{jr}(V) \cup \phi_{kr}(V)|^3 \geq \frac{1}{6^2}\left(\frac{n(n-1)(n-2)R}{P}\right)^3,$$
$$\frac{n(n-1)(n-2)}{6P} \leq |\phi_{ijk}(V)| \leq \frac{n(n-1)(n-2)}{6},$$
$$\frac{nR}{P} \leq |\phi_{ir}(V) \cup \phi_{jr}(V) \cup \phi_{kr}(V)| \leq nR,$$
where $P,n$, and $R$ are all positive integers greater than or equal to 1.
\heading{Minimum communication (bandwidth) cost}
\emph{Communication lower bound} = $|\phi_{ijk}(V)|^* + 2|\phi_{ir}(V) \cup \phi_{jr}(V) \cup \phi_{kr}(V)|^* -W$, where $W$ denotes the data owned by the processor.

We assume that there is one copy of input and output data at the beginning and the end. There exists a processor who has at most $1/Pth$ of overall data. Therefore, $W=\frac{n(n-1)(n-2)}{6P} + 2\frac{nR}{P}$.

Usually $R$ is not very large. For $R \leq \frac{(n-1)(n-2)}{2^3}$ and small $P$, \emph{communication lower bound} $= 2\left(\frac{n(n-1)(n-2)}{P}\right)^{\frac13}R - 2\frac{n}{P}R\text.$

%For the typical case where $R$ is not very large, ($R \leq \frac{(n-1)(n-2)}{2^3}$), and small $P$, \emph{communication lower bound} $= 2\left(\frac{n(n-1)(n-2)}{P}\right)^{1/3}R - 2\frac{nR}{P}\text.$
\end{block}
\begin{block}{Communication optimal algorithm for symmetric MTTKRP}
\textcolor{brown}{Approach}: express computations in terms of tetrahedron blocks

\begin{itemize}
	\item construct equal size sets such that each triplet is covered and appears in exactly one set
	\item assign diagonal blocks such that the remaining load is almost balanced and no new entries of the symmetric tensor are required
\end{itemize}

\vspace*{-0.35cm}
\heading{An example for $10\times 10 \times 10$ blocks of a symmetric tensor on $30$ processors.}
%	\begin{tabular}{|c|c|c|c|c|}
%	\hline
%	$p$ & $R_p$ & $N_p$ & $D_p$ \\
%	\hline
%	1 & \{1,2,3,7\} & \{(2,2,1), (2,1,1), (7,2,2)\} & \{(1,1,1)\} \\
%	2 & \{1,2,4,5\} & \{(4,4,1), (4,1,1), (5,1,1)\} & \{(2,2,2)\} \\
%	3 & \{1,2,6,10\} & \{(6,6,1),(10,10,2), (6,1,1)\} & \{(6,6,6)\} \\
%	4 & \{1,2,8,9\} & \{(8,8,1),(9,9,8), (8,1,1)\} & \{(8,8,8)\} \\
%	5 & \{1,3,4,10\} & \{(10,10,1),(10,10,3), (10,1,1)\} & \{(3,3,3)\} \\
%	6 & \{1,3,5,8\} & \{(3,3,1),(8,8,5), (3,1,1)\} & \{(5,5,5)\} \\
%	7 & \{1,3,6,9\} & \{(9,9,1),(9,9,3), (9,1,1)\} & \{(9,9,9)\} \\
%	8 & \{1,4,6,8\} & \{(6,6,4),(8,8,6), (6,4,4)\} & \{(4,4,4)\} \\
%	9 & \{1,4,7,9\} & \{(7,7,1), (9,9,4), (7,1,1))\} & \{(7,7,7)\} \\
%	10 & \{1,5,6,7\} & \{(5,5,1), (7,7,6), (7,6,6)\} & \{\} \\
%	11 & \{1,5,9,10\} & \{(9,9,5),(10,10,9), (9,5,5)\} & \{(10,10,10)\}  \\
%	12 & \{1,7,8,10\} & \{(8,8,7),(10,10,8), (10,8,8)\} & \{\} \\
%	13 & \{2,3,4,8\} & \{(3,3,2), (3,2,2),(4,2,2)\} & \{\}  \\
%	14 & \{2,3,5,6\} & \{(5,5,2), (5,2,2),(6,5,5)\} & \{\}  \\
%	15 & \{2,3,9,10\} & \{(9,9,2), (9,2,2),(10,2,2)\} & \{\}  \\
%	16 & \{2,4,6,9\} & \{(4,4,2), (9,9,6), (9,6,6)\} & \{\}  \\
%	17 & \{2,4,7,10\} & \{(7,7,2),(10,10,4), (10,4,4)\} & \{\}  \\
%	18 & \{2,5,7,9\} & \{(7,7,5),(9,9,7), (7,5,5)\} & \{\} \\
%	19 & \{2,5,8,10\} & \{(8,8,2), (8,2,2),(10,5,5)\} & \{\}  \\
%	20 & \{2,6,7,8\} & \{(6,6,2), (6,2,2),(8,6,6)\} & \{\}  \\
%	21 & \{3,4,5,9\} & \{(4,4,3), (4,3,3),(9,4,4)\} & \{\}  \\
%	22 & \{3,4,6,7\} & \{(6,6,3), (6,3,3),(7,3,3)\} & \{\}  \\
%	23 & \{3,5,7,10\} & \{(5,5,3), (5,3,3),(10,3,3)\} & \{\}  \\
%	24 & \{3,6,8,10\} & \{(8,8,3),(10,10,6), (8,3,3)\} & \{\}  \\
%	25 & \{3,7,8,9\} & \{(7,7,3), (9,7,7),(9,3,3)\} & \{\}  \\
%	26 & \{4,5,6,10\} & \{(5,5,4), (5,4,4),(10,10,5)\} & \{\}  \\
%	27 & \{4,5,7,8\} & \{(7,7,4), (7,4,4),(8,7,7)\} & \{\}  \\
%	28 & \{4,8,9,10\} & \{(8,8,4), (8,4,4),(10,9,9)\} & \{\}  \\
%	29 & \{5,6,8,9\} & \{(6,6,5), (8,5,5) (9,8,8)\} & \{\} \\
%	30 & \{6,7,9,10\} & \{(10,6,6),(10,10,7), (10,7,7)\} & \{\}  \\
%	\hline
%\end{tabular}

\begin{center}
	\scalebox{0.965}{
	\begin{tabular}{|c|c|c|c|c|}
	\hline
	$p$ & $R_p$ & $N_p$ & $D_p$ \\
		\hline
		1  & $\{1,2,3,4\}$ & $\{(2,2,1),(2,1,1),(3,3,1)\}$ & $\{(1,1,1)\}$ \\
		2  & $\{1,2,5,6\}$ & $\{(6,6,1),(6,1,1),(6,2,2)\}$ & $\{(2,2,2)\}$ \\
		3  & $\{1,2,7,8\}$ & $\{(7,7,1),(7,1,1),(8,8,7)\}$ & $\{(7,7,7)\}$ \\
		4  & $\{1,2,9,10\}$ & $\{(10,10,1),(10,1,1),(9,2,2)\}$ & $\{(9,9,9)\}$ \\
		5  & $\{1,3,5,7\}$ & $\{(3,1,1),(5,5,3),(5,3,3)\}$ & $\{(3,3,3)\}$ \\
		6  & $\{1,3,6,9\}$ & $\{(9,9,1),(9,1,1),(9,9,6)\}$ & $\{(6,6,6)\}$ \\
		7  & $\{1,3,8,10\}$ & $\{(8,8,1),(8,1,1),(10,10,8)\}$ & $\{(10,10,10)\}$ \\
		8  & $\{1,4,5,10\}$ & $\{(4,4,1),(4,1,1),(5,4,4)\}$ & $\{(5,5,5)\}$ \\
		9  & $\{1,4,6,8\}$ & $\{(6,6,4),(6,4,4),(8,8,4)\}$ & $\{(8,8,8)\}$ \\
		10 & $\{1,4,7,9\}$ & $\{(7,7,4),(7,4,4),(9,9,7)\}$ & $\{(4,4,4)\}$ \\
		11 & $\{1,5,8,9\}$ & $\{(5,5,1),(5,1,1),(9,5,5)\}$ & $\{\}$ \\
		12 & $\{1,6,7,10\}$ & $\{(10,10,6),(10,6,6),(10,10,7)\}$ & $\{\}$ \\
		13 & $\{2,3,5,8\}$ & $\{(3,3,2),(3,2,2),(8,5,5)\}$ & $\{\}$ \\
		14 & $\{2,3,6,10\}$ & $\{(10,10,3),(10,3,3),(6,6,2)\}$ & $\{\}$ \\
		15 & $\{2,3,7,9\}$ & $\{(9,9,3),(9,3,3),(9,7,7)\}$ & $\{\}$ \\
		16 & $\{2,4,5,9\}$ & $\{(4,4,2),(4,2,2),(9,9,2)\}$ & $\{\}$ \\
		17 & $\{2,4,6,7\}$ & $\{(7,7,2),(7,2,2),(7,6,6)\}$ & $\{\}$ \\
		18 & $\{2,4,8,10\}$ & $\{(8,8,2),(8,2,2),(10,10,2)\}$ & $\{\}$ \\
		19 & $\{2,5,7,10\}$ & $\{(5,5,2),(5,2,2),(10,2,2)\}$ & $\{\}$ \\
		20 & $\{2,6,8,9\}$ & $\{(8,6,6),(9,6,6),(9,8,8)\}$ & $\{\}$ \\
		21 & $\{3,4,5,6\}$ & $\{(4,4,3),(4,3,3),(5,5,4)\}$ & $\{\}$ \\
		22 & $\{3,4,7,10\}$ & $\{(7,7,3),(7,3,3),(10,10,4)\}$ & $\{\}$ \\
		23 & $\{3,4,8,9\}$ & $\{(8,8,3),(8,3,3),(8,4,4)\}$ & $\{\}$ \\
		24 & $\{3,5,9,10\}$ & $\{(9,9,5),(10,10,9),(10,9,9)\}$ & $\{\}$ \\
		25 & $\{3,6,7,8\}$ & $\{(6,6,3),(6,3,3),(8,8,6)\}$ & $\{\}$ \\
		26 & $\{4,5,7,8\}$ & $\{(7,7,5),(7,5,5),(8,8,5)\}$ & $\{\}$ \\
		27 & $\{4,6,9,10\}$ & $\{(9,9,4),(9,4,4),(10,4,4)\}$ & $\{\}$ \\
		28 & $\{5,6,7,9\}$ & $\{(6,6,5),(6,5,5),(7,7,6)\}$ & $\{\}$ \\
		29 & $\{5,6,8,10\}$ & $\{(10,10,5),(10,5,5),(10,8,8)\}$ & $\{\}$ \\
		30 & $\{7,8,9,10\}$ & $\{(8,7,7),(10,7,7),(9,9,8)\}$ & $\{\}$ \\
		\hline
	\end{tabular}
}
\end{center}

$R_p$: computations associated with each triplet of this set are performed on processor $p$\\
$N_p \text{ and } D_p$: computations associated with diagonal blocks on processor $p$

%\vspace*{-0.21cm}
\heading{Work performed by processor $5$}
%As $R_7= \{1,3,6,9\},  N_7 = \{(9,9,1),(9,9,3), (9,1,1)\} \text{ and } D_7 \{(9,9,9)\}$. THis processor will perform all computations associated with $(6,3,1), (9,3,1), (9,6,1), (9,6,3)$ off-diagonal tiles and $(9,9,1),(9,9,3), (9,1,1), (9,9,9)$ diagonal tiles.  

As $R_5= \{1,3,5,7\},$  $N_5 = \{(3,1,1),(5,5,3),(5,3,3)\} \text{ and } D_5 = \{(3,3,3)\}$. Processor $5$ performs all computations associated with $(5,3,1), (7,3,1), (7,5,1), (7,5,3)$ off-diagonal blocks and $(3,1,1),(5,5,3),(5,3,3), (3,3,3)$ diagonal blocks.  

%$\bullet$  A tetrahedron block partition for $n \times n \times $ blocks symmetric tensor exists on $P$ processors when $n=q^2+1$ and $P=q(q^2+1)$ for some prime $q$ (by the properties of Steiner systems)
\vspace*{-0.25cm}
\begin{itemize}
	\item A tetrahedron block partition for $n \times n \times n$ blocks symmetric tensor exists on $P$ processors when $n=q^2+1$ and $P=q(q^2+1)$ for some prime $q$ (by the properties of Steiner systems) 
	%	\item Existence of the bound
\end{itemize}

\end{block}


%  \begin{block}{References}
%
%    \nocite{*}
%    \footnotesize{\bibliographystyle{plain}\bibliography{poster}}
%
%  \end{block}

\end{column}

\separatorcolumn
\end{columns}
\end{frame}

\end{document}
